\scanpage{267}
\mspage{11}
\margin{n° 227}
$n$ sur $Q$. Mais $z-aw$ est dans $L(D+Q)$, donc dans $L(D)$. Ceci contredit l'indépendence linéaire de $z$ et $w$ mod $L(D)$.

Supposons maintenant le théorème démontré pour toutes les variétés de dimension $<r$. On va procéder par induction. Faisons d'abord quelques remarques. Soit $z$ une fonction sur $V$, et $W$ un diviseur premier de $V$. Alors $W$ n'est pas nécessairement normale. Soit $P'$ un point générique de $W$. Alors la spécialisation $P\longrightarrow P'$ induit une place du corps $k(P)$. Si la fonction $z$ n'a pas de pôle sur $W$, elle est dans l'anneau de valuation de cette place et est définie sur $W$. Son image $z'$ dans le corps $k(P')$ définit une fonction sur $W$. La normalisation $W^*$ de $W$ ayant le même corps de fonction que $W$, on peut dire que si $z$ n'a pas de pôle sur $W$, alors $z$ induit sur $W$ ou sur $W^*$ une fonction $z'$. C'est de cette fonction qu'il s'agira dans le lemme suivant.

LEMME. -- \uline{Soit} $D$ \uline{un diviseur sur} $V$, \uline{et} $W$ \uline{un diviseur premier de} $V$, \uline{non contenu dans les composantes polaires de} $D$. \uline{Alors il existe un diviseur} $D^*$ \uline{sur la normalisation} $W^*$ \uline{de} $W$ \uline{tel que pour une fonction} $z\in L(D)$, \uline{la fonction induite} $z'$ \uline{sur} $W^*$ \uline{soit dans} $L(D^*)$.

\uline{Dém}. Soit $V_i$, $i=1,\ldots,g$ les représentants de $V$. On peut supposer que tous les $V_i$ sont des variétés affines. D'après un théorème d'approximation élémentaire, il existe une fonction $f_i\in K$ telle que pour $z\in L(D)$, la fonction $f_i z$ n'ait pas de pôle ayant un représentant sur $V_i$, et telle que $f_i$ soit une unité sur $W$.

Si $W$ a un représentant sur $V_i$, alors les fonctions $f_i$, $z$, et $f_i z$ sont définies sur $W$, et induisent des fonctions $f'_i$, $z'$ et $f'_i z'$ sur $W$. De plus, $\unsure{f'_j|f'_i}\neq0$. Alors, d'après la prop.5 ($V$ étant supposée

\mspage{266}
